%! Author = Saeid Yazdani
%! Date = 12/17/2023


\section{The Need for Modern Test Equipment}
\label{sec:need-modern-test-equipment}

The growth in \ac{IC} technology is driven by the need for more
functions per device while simultaneously reducing power consumption to comply
with environmentally friendly directives and extend the operational time of battery
powered applications.
This demand pushes the evolution of \ac{IC} technology almost on a daily basis.

Logic and memory cells have traditionally driven reductions in technology size \autocite{8776580}, 
and the increasing integration density of digital circuits raises the complexity of digital logic 
testing. At the same time, high-speed \ac{I/O} interfaces such as \ac{PCIe}, \ac{USB}, and \ac{HDMI} 
are now common in high-volume consumer products. These interfaces introduce challenges more typical 
of analog and RF testing, including signal integrity, timing, and jitter, which significantly impact 
test coverage, cost, and overall testing strategy.

With each new technology, novel failure mechanisms emerge, making advanced
silicon debugging and yield learning critical during the ramp-up of
processes.
Additionally, emerging trends in 3D packaging technologies are pushing the
limits of testing at the probe stage \autocite{8480435}.

\ac{EV} development has introduced the demand for smart components with high power ratings. 
Initially, discrete switching components needed \SI{70}{\volt} to \SI{100}{\volt} maximum drain 
voltage for power rails of \SI{12}{\volt} to \SI{48}{\volt}, with currents typically below 
\SI{50}{\ampere}. With the introduction of \ac{EV}s, switching components capable of withstanding 
up to \SI{1}{\kilo\volt} and currents in the range of \SI{100}{\ampere} to \SI{1000}{\ampere} 
became necessary, as motor power reaches several hundred kilowatts. This trend initiated the demand 
for high-power components, and today state-of-the-art \ac{ATE} features both low-voltage and 
high-voltage sections designed to handle such high currents safely and reliably.

When it comes to testing and qualification of components, there is a broad
range of properties associated with \ac{DUT}s that need testing and
qualification:

\begin{itemize}[topsep=8pt,itemsep=4pt,partopsep=4pt, parsep=4pt]
    \item \textbf{\emph{Type:}} Is the \ac{DUT} a discrete component, an integrated solution, or a MEMS device?
    \item \textbf{\emph{Power rating:}} Operating voltage and current ranges, including supply voltages.
    \item \textbf{\emph{Pin Count:}} Number of \ac{I/O}, control, and power pins.
    \item \textbf{\emph{Function:}} Functions of the \ac{DUT}: analog, digital, mixed-signal, or MEMS.
    \item \textbf{\emph{Frequency:}} Operating frequency or bandwidth specifications for signal processing or MEMS actuation.
\end{itemize}

The vast number of different components and, in the case of \ac{SoC}s, different
functional blocks, would require an \ac{ATE} that can cover all possible
scenarios.
This raises the question for an most economical yet reliable solution
for testing and qualification.

Today, a single \ac{ATE} is responsible for performing
all the required tests on a given \ac{DUT}.
This test strategy is mandatory when the functional blocks of a \ac{DUT} can
not be tested independently.
For example, a mixed-signal \ac{IC} may require simultaneous stimulation of its
analog block to check for the correctness of its digital block's output.

The wide range of technological and functional variations among components makes it challenging 
for a single \ac{ATE} to cover all possible test scenarios, including different power ratings 
and operating frequencies. Modern mixed-signal devices and \ac{SoC}s, however, can often be 
tested on a single tester, since most or all functional blocks can be active simultaneously.

\subsection{How \acs{DfT} Influences \acs{ATE} requirements}
\label{subsec:properties-of-modern-ate}

Due to advances in the equipment
industry, modern \ac{ATE}s exist that cover all
test requirements and scenarios for state-of-the-art \acp{DUT}.

Such \ac{ATE}s are flexible and modular, allowing for on-demand
removal and insertion of modules to provide test facilities.
Such modules include:

\begin{itemize}[topsep=8pt,itemsep=4pt,partopsep=4pt, parsep=4pt]
    \item Analog signal processing with \ac{DSP} instrumentation.
    \item Analog power processing with \ac{SMU}s in
    \SI{}{\kilo\volt} or \SI{}{\kilo\ampere} range.
    \item Digital core testing, SCAN or \ac{BIST}.
    \item Memory test with algorithmic pattern generator hardware\footnote{
        Modern SoCs incorporate Memory BIST that is accessed via SCAN paths.}.
    \item \ac{RF} testing in \SI{}{\giga\hertz} range for modern communication.
\end{itemize}

A modern \ac{ATE} needs to provide a set of key features:

\begin{itemize}[topsep=8pt,itemsep=4pt,partopsep=4pt, parsep=4pt]

    \item \textbf{\emph{High-Speed:}} \ac{ATE}'s rate of stimulation
    should be high enough to test all functions of \acp{DUT}.

    \item \textbf{\emph{Precise and accurate:}} Confidence of PASS/FAIL results
    are highly dependent on the precision and accuracy of generated
    and measured signals in terms of voltage, current and timing.

    \item \textbf{\emph{Parallelism:}} \acp{ATE} are able to
    test multiple \acp{DUT} at the same time through multi-site capabilities.

    \item \textbf{\emph{Scalability:}} Modern \acp{ATE}
    keep up with new technologies by allowing hardware and
    software updatability without the need for a total overhaul or going
    obsolete.

\end{itemize}



